%% toc: P2 9 9.5,9.6,9.7
\section{Proofs for the fourteen syllogism rules}
\label{sec:annex-p02-w02-p02-002-fourteen}

Throughout, \(\mathfrak A\) is a structure with nonempty universe \(A\), as
in Definition~\ref{def:p02-syllogistic-semantics}, and term letters are read
as atoms by Convention~\ref{def:p02-w02-p02-002-reading}. For atoms \(z,x\)
recall the four truth conditions
\[
 \mathfrak A\models\forall(z,x)\iff z^{\mathfrak A}\subseteq x^{\mathfrak A},
 \qquad
 \mathfrak A\models\forall(z,\bar x)\iff z^{\mathfrak A}\cap x^{\mathfrak A}=\varnothing,
\]
\[
 \mathfrak A\models\exists(z,x)\iff z^{\mathfrak A}\cap x^{\mathfrak A}\ne\varnothing,
 \qquad
 \mathfrak A\models\exists(z,\bar x)\iff z^{\mathfrak A}\setminus x^{\mathfrak A}\ne\varnothing,
\]
which follow from Definition~\ref{def:p02-syllogistic-semantics} because
\((\bar x)^{\mathfrak A}=A\setminus x^{\mathfrak A}\). Non-derivability
claims use soundness of \(\mathsf S\)
(Theorem~\ref{thm:p02-syllogistic-s-completeness}): if a structure satisfies
every premise and falsifies the conclusion, there is no derivation. In each
countermodel, atoms not mentioned are interpreted by \(\varnothing\).

\begin{annexitem}{annex:p02-w02-p02-002-conversions}{thm:p02-w02-p02-002-conversions}{Conversions}
\begin{proof}
\emph{Item 1.} Definition~\ref{def:p02-syllogistic-identifications}
identifies \(\forall(l,m)\) with \(\forall(\bar m,\bar l)\). With \(l=b\) and
\(m=\bar a\), and since \(\overline{\bar a}=a\), this identifies
\(\forall(b,\bar a)\) with \(\forall(a,\bar b)\). By
Definition~\ref{def:p02-syllogistic-derivation} a derivation may use the
declared identifications, so the one-node tree whose leaf is the premise
\(\forall(b,\bar a)\) is a derivation of \(\forall(a,\bar b)\), and
conversely.

\emph{Item 2.} The same definition identifies \(\exists(b,a)\) with
\(\exists(a,b)\); the argument of item~1 applies.

\emph{Item 3.} Let \(A=\{c\}\), \(a^{\mathfrak A}=\{c\}\) and
\(b^{\mathfrak A}=\varnothing\). Then \(b^{\mathfrak A}\subseteq
a^{\mathfrak A}\), so \(\forall(b,a)\) is true, and
\(a^{\mathfrak A}\cap b^{\mathfrak A}=\varnothing\), so \(\exists(a,b)\) is
false. By soundness, \(\forall(b,a)\nvdash_{\mathsf S}\exists(a,b)\).
With \(\exists(b,b)\) added, apply D1,
\(\forall(q,l),\exists(p,q)\mathbin{/}\exists(p,l)\), with \(q:=b\),
\(l:=a\) and \(p:=b\): the antecedents are \(\forall(b,a)\) and
\(\exists(b,b)\), and the conclusion is \(\exists(b,a)\), which is
identified with \(\exists(a,b)\).
\end{proof}
\end{annexitem}

\begin{annexitem}{annex:p02-w02-p02-002-twelve}{thm:p02-w02-p02-002-twelve}{Twelve rules}
\begin{proof}
For each rule we name the rule of \(\mathsf S\) and the values of its
metavariables. The rules are, from the sheet for \(\mathsf S\) in Chapter~\ref{sec:p02-syllogistic-inference},

\begin{itemize}
\item B: \(\forall(p,q),\forall(q,l)\mathbin{/}\forall(p,l)\);
\item D1: \(\forall(q,l),\exists(p,q)\mathbin{/}\exists(p,l)\);
\item D2: \(\exists(p,l),\forall(p,q)\mathbin{/}\exists(q,l)\);
\item D3: \(\forall(q,\bar l),\exists(p,l)\mathbin{/}\exists(p,\bar q)\),
\end{itemize}
where \(p,q\) are atoms and \(l\) is a literal. In an assignment
such as \(p:=s\), the left side is the metavariable of the rule and the right
side the atom assigned to it; in the third figure the atom \(p\) may also
occur on a right side.

\emph{SR1 Barbara.} B with \(p:=c\), \(q:=b\), \(l:=a\) has antecedents
\(\forall(c,b)\), \(\forall(b,a)\) and conclusion \(\forall(c,a)\).

\emph{SR2 Celarent.} B with \(p:=c\), \(q:=b\), \(l:=\bar a\) has antecedents
\(\forall(c,b)\), \(\forall(b,\bar a)\) and conclusion \(\forall(c,\bar a)\).

\emph{SR3 Darii.} D1 with \(q:=b\), \(l:=a\), \(p:=c\) has antecedents
\(\forall(b,a)\), \(\exists(c,b)\) and conclusion \(\exists(c,a)\).

\emph{SR4 Ferio.} D1 with \(q:=b\), \(l:=\bar a\), \(p:=c\) has antecedents
\(\forall(b,\bar a)\), \(\exists(c,b)\) and conclusion \(\exists(c,\bar a)\).

\emph{SR5 Cesare.} The premise \(\forall(n,\bar m)\) is identified with
\(\forall(m,\bar n)\) (Lemma~\ref{thm:p02-w02-p02-002-conversions}, item~1).
B with \(p:=x\), \(q:=m\), \(l:=\bar n\) has antecedents \(\forall(x,m)\),
\(\forall(m,\bar n)\) and conclusion \(\forall(x,\bar n)\).

\emph{SR6 Camestres.} The premise \(\forall(x,\bar m)\) is identified with
\(\forall(m,\bar x)\). B with \(p:=n\), \(q:=m\), \(l:=\bar x\) has
antecedents \(\forall(n,m)\), \(\forall(m,\bar x)\) and conclusion
\(\forall(n,\bar x)\), which is identified with \(\forall(x,\bar n)\).

\emph{SR7 Festino.} The premise \(\forall(n,\bar m)\) is identified with
\(\forall(m,\bar n)\). D1 with \(q:=m\), \(l:=\bar n\), \(p:=x\) has
antecedents \(\forall(m,\bar n)\), \(\exists(x,m)\) and conclusion
\(\exists(x,\bar n)\).

\emph{SR8 Baroco.} D3 with \(q:=n\), \(l:=\bar m\), \(p:=x\). Since
\(\overline{\bar m}=m\), its antecedents are \(\forall(n,m)\) and
\(\exists(x,\bar m)\), and its conclusion is \(\exists(x,\bar n)\).

\emph{SR11 Disamis.} D2 with \(p:=s\), \(l:=p\), \(q:=r\) has antecedents
\(\exists(s,p)\), \(\forall(s,r)\) and conclusion \(\exists(r,p)\).

\emph{SR12 Datisi.} The premise \(\exists(s,r)\) is identified with
\(\exists(r,s)\). D1 with \(q:=s\), \(l:=p\), \(p:=r\) has antecedents
\(\forall(s,p)\), \(\exists(r,s)\) and conclusion \(\exists(r,p)\).

\emph{SR13 Bocardo.} D2 with \(p:=s\), \(l:=\bar p\), \(q:=r\) has
antecedents \(\exists(s,\bar p)\), \(\forall(s,r)\) and conclusion
\(\exists(r,\bar p)\).

\emph{SR14 Ferison.} The premise \(\exists(s,r)\) is identified with
\(\exists(r,s)\). D1 with \(q:=s\), \(l:=\bar p\), \(p:=r\) has antecedents
\(\forall(s,\bar p)\), \(\exists(r,s)\) and conclusion \(\exists(r,\bar p)\).

In each case the values of the atom metavariables are atoms and the value
of \(l\) is a literal, as the side conditions of the sheet require, and every
formula displayed is an \(\mathcal S\)-formula, because its first literal is
an atom (Definition~\ref{def:p02-syllogistic-language}). Each derivation is a
tree with the two premises as leaves and one rule node, read up to the
stated identifications. The semantic consequences follow by soundness of
\(\mathsf S\) (Theorem~\ref{thm:p02-syllogistic-s-completeness}), which
places no condition on the emptiness of any extension.
\end{proof}
\end{annexitem}

\begin{annexitem}{annex:p02-w02-p02-002-third-figure}{thm:p02-w02-p02-002-third-figure}{Darapti and Felapton}
\begin{proof}
Write \(S=s^{\mathfrak A}\), \(R=r^{\mathfrak A}\), \(P=p^{\mathfrak A}\).

\emph{Item 1, if.} Suppose \(\exists(s,s)\in E\). For \(k=9\), D1 with
\(q:=s\), \(l:=p\), \(p:=s\) applied to \(\forall(s,p)\) and \(\exists(s,s)\)
gives \(\exists(s,p)\); D2 with \(p:=s\), \(l:=p\), \(q:=r\) applied to
\(\exists(s,p)\) and \(\forall(s,r)\) gives \(\exists(r,p)=\theta_9\).
For \(k=10\), the same two steps with \(l:=\bar p\) give
\(\exists(s,\bar p)\) from \(\forall(s,\bar p)\) and \(\exists(s,s)\), and
then \(\exists(r,\bar p)=\theta_{10}\) from \(\exists(s,\bar p)\) and
\(\forall(s,r)\). The derivation uses only premises in
\(\Gamma_k\cup\{\exists(s,s)\}\subseteq\Gamma_k\cup E\), and a derivation
from a set is a derivation from every superset, since its leaves remain
premises.

\emph{Item 1, only if.} Suppose \(\exists(s,s)\notin E\), so
\(E\subseteq\{\exists(p,p),\exists(r,r)\}\). For \(k=9\) take \(A=\{c,d\}\)
with \(c\ne d\), \(S=\varnothing\), \(P=\{c\}\), \(R=\{d\}\). Both members of
\(\Gamma_9\) have the empty subject \(S\) and are true.
\(\exists(p,p)\) and \(\exists(r,r)\) are true because \(P\) and \(R\) are
nonempty. \(R\cap P=\varnothing\), so \(\theta_9\) is false. For \(k=10\)
take \(A=\{c\}\), \(S=\varnothing\), \(P=R=\{c\}\). Both members of
\(\Gamma_{10}\) are true because \(S\) is empty, \(\exists(p,p)\) and
\(\exists(r,r)\) are true, and \(R\setminus P=\varnothing\), so
\(\theta_{10}\) is false. In both cases the structure satisfies
\(\Gamma_k\cup E\) for every such \(E\) and falsifies \(\theta_k\);
by soundness \(\Gamma_k\cup E\nvdash_{\mathsf S}\theta_k\). The case
\(E=\varnothing\) gives \(\Gamma_k\nvdash_{\mathsf S}\theta_k\).

\emph{Item 2.} If \(\Theta'\subseteq\Theta''\) and a structure satisfies
\(\Theta''\) while falsifying \(\theta\), it satisfies \(\Theta'\); so
\(\Theta''\not\models\theta\) implies \(\Theta'\not\models\theta\). Every
proper subset of the three-element set
\(\Gamma_k\cup\{\exists(s,s)\}\) is contained in one of its three
two-element subsets, so it suffices to refute those.
\begin{itemize}
\item Without \(\exists(s,s)\): the set is \(\Gamma_k\), refuted in item~1.
\item Without \(\forall(s,r)\). For \(k=9\), take \(A=\{c\}\), \(S=P=\{c\}\),
 \(R=\varnothing\): \(\forall(s,p)\) and \(\exists(s,s)\) are true and
 \(R\cap P=\varnothing\). For \(k=10\), take \(A=\{c\}\), \(S=\{c\}\),
 \(P=R=\varnothing\): \(S\cap P=\varnothing\) makes \(\forall(s,\bar p)\)
 true, \(\exists(s,s)\) is true, and \(R\setminus P=\varnothing\).
\item Without the premise containing \(p\). For \(k=9\), take \(A=\{c\}\),
 \(S=R=\{c\}\), \(P=\varnothing\): \(\forall(s,r)\) and \(\exists(s,s)\) are
 true and \(R\cap P=\varnothing\). For \(k=10\), take \(A=\{c\}\),
 \(S=R=P=\{c\}\): \(\forall(s,r)\) and \(\exists(s,s)\) are true and
 \(R\setminus P=\varnothing\).
\end{itemize}
In each case \(\theta_k\) is false in a structure satisfying the two
remaining premises.

\emph{Item 3.} By Definition~\ref{def:p02-syllogistic-identifications} the
only identifications are \(\exists(l,m)\sim\exists(m,l)\) and
\(\forall(l,m)\sim\forall(\bar m,\bar l)\), so an existential formula is
identified only with existential formulas. The only existential member of
\(\Gamma_k\cup\{\exists(s,s)\}\) is \(\exists(s,s)\), whose identification
class is \(\{\exists(s,s)\}\). The class of \(\theta_9\) is
\(\{\exists(r,p),\exists(p,r)\}\) and that of \(\theta_{10}\) is
\(\{\exists(r,\bar p),\exists(\bar p,r)\}\). Since \(s\), \(r\) and \(p\) are
distinct atoms, neither class contains \(\exists(s,s)\).
\end{proof}
\end{annexitem}

\begin{annexitem}{annex:p02-w02-p02-002-committing}{thm:p02-w02-p02-002-committing}{Existential import}
\begin{proof}
For a sentence \(XyZ\) put \(Z^\ast=z^{\mathfrak A}\) and
\(X^\ast=x^{\mathfrak A}\). By Definition~\ref{def:p02-w02-p02-001-patterns},
\(\mathfrak A\models_{\mathrm B}XyZ\) holds iff the condition
\(\alpha,\varepsilon,\iota,\omega\) for \(y=a,e,i,o\), respectively, holds of
\((Z^\ast,X^\ast)\), with these conditions as in
Definition~\ref{def:p13-boger-guarded-bridge}. Unfolding them and comparing
with the four truth conditions at the head of this section gives three facts.
\begin{enumerate}
\item[(F1)] \(\mathfrak A\models_{\mathrm B}XaZ\) iff \(Z^\ast\ne\varnothing\)
 and \(\mathfrak A\models\sigma(XaZ)\), since
 \(\alpha(Z^\ast,X^\ast)\) is \(Z^\ast\ne\varnothing\land Z^\ast\subseteq X^\ast\).
\item[(F2)] \(\mathfrak A\models_{\mathrm B}XiZ\) iff
 \(\mathfrak A\models\sigma(XiZ)\), since \(\iota(Z^\ast,X^\ast)\) is
 \(Z^\ast\cap X^\ast\ne\varnothing\); and
 \(\mathfrak A\models_{\mathrm B}XeZ\) iff \(\mathfrak A\models\sigma(XeZ)\),
 since \(\varepsilon\) is the negation of \(\iota\), that is,
 \(Z^\ast\cap X^\ast=\varnothing\).
\item[(F3)] \(\mathfrak A\models_{\mathrm B}XoZ\) iff \(Z^\ast=\varnothing\)
 or \(\mathfrak A\models\sigma(XoZ)\). Indeed \(\omega\) is the negation of
 \(\alpha\), which holds iff \(Z^\ast=\varnothing\) or
 \(Z^\ast\not\subseteq X^\ast\), and \(Z^\ast\not\subseteq X^\ast\) iff
 \(Z^\ast\setminus X^\ast\ne\varnothing\).
\end{enumerate}
In particular \(\models_{\mathrm B}\) implies \(\sigma\)-truth for a-,
e- and i-sentences, and \(\sigma\)-truth implies \(\models_{\mathrm B}\) for e-,
i- and o-sentences. Let \(\mathfrak A\) satisfy both premises of SR\(k\)
under \(\models_{\mathrm B}\).

\emph{Rules SR2--SR7, SR11, SR12 and SR14.} None of them has an o-premise,
and each has an e-, i- or o-conclusion. Their premises are a-, e- or i-sentences, so by (F1) and (F2)
\(\mathfrak A\) satisfies their \(\sigma\)-readings. By
Theorem~\ref{thm:p02-w02-p02-002-twelve}, \(\mathfrak A\models\sigma(\psi)\)
for the conclusion \(\psi\). As \(\psi\) is an e-, i- or o-sentence,
\(\mathfrak A\models_{\mathrm B}\psi\) by (F2) or (F3).

\emph{SR1 Barbara.} From \(BaC\), by (F1), \(c^{\mathfrak A}\ne\varnothing\)
and \(\mathfrak A\models\forall(c,b)\); from \(AaB\),
\(\mathfrak A\models\forall(b,a)\). By
Theorem~\ref{thm:p02-w02-p02-002-twelve}, \(\mathfrak A\models\forall(c,a)\).
With \(c^{\mathfrak A}\ne\varnothing\), (F1) gives
\(\mathfrak A\models_{\mathrm B}AaC\).

\emph{SR9 Darapti and SR10 Felapton.} The premise \(RaS\) gives, by (F1),
\(s^{\mathfrak A}\ne\varnothing\), that is \(\mathfrak A\models\exists(s,s)\),
and \(\mathfrak A\models\forall(s,r)\). The other premise, \(PaS\) or
\(PeS\), gives \(\mathfrak A\models\sigma(PaS)\) or
\(\mathfrak A\models\sigma(PeS)\) by (F1) or (F2). Thus \(\mathfrak A\)
satisfies \(\Gamma_k\cup\{\exists(s,s)\}\). By
Theorem~\ref{thm:p02-w02-p02-002-third-figure}, item~1, and soundness,
\(\mathfrak A\models\theta_k\). For \(k=9\), \(\theta_9=\sigma(PiR)\) and
(F2) gives \(\mathfrak A\models_{\mathrm B}PiR\); for \(k=10\),
\(\theta_{10}=\sigma(PoR)\) and (F3) gives \(\mathfrak A\models_{\mathrm B}PoR\).
This also proves the second sentence of the theorem.

\emph{SR8 Baroco.} If \(x^{\mathfrak A}=\varnothing\), then the conclusion
\(NoX\) has an empty subject and holds under \(\models_{\mathrm B}\) by (F3).
Otherwise the premise \(MoX\) with nonempty subject gives
\(\mathfrak A\models\sigma(MoX)\) by (F3), and \(MaN\) gives
\(\mathfrak A\models\sigma(MaN)\) by (F1). By
Theorem~\ref{thm:p02-w02-p02-002-twelve},
\(\mathfrak A\models\sigma(NoX)\), and (F3) gives
\(\mathfrak A\models_{\mathrm B}NoX\).

\emph{SR13 Bocardo.} The premise \(RaS\) gives \(s^{\mathfrak A}\ne\varnothing\)
and \(\mathfrak A\models\sigma(RaS)\) by (F1). The premise \(PoS\) has the
nonempty subject \(s^{\mathfrak A}\), so (F3) gives
\(\mathfrak A\models\sigma(PoS)\). By
Theorem~\ref{thm:p02-w02-p02-002-twelve},
\(\mathfrak A\models\sigma(PoR)\), and (F3) gives
\(\mathfrak A\models_{\mathrm B}PoR\).

The fourteen cases are exhausted. For the last sentence of the theorem:
the twelve rules other than SR9 and SR10 hold in \(\mathsf S\)
(Theorem~\ref{thm:p02-w02-p02-002-twelve}) and under \(\models_{\mathrm B}\);
SR9 and SR10 hold under \(\models_{\mathrm B}\) but not in \(\mathsf S\)
(Theorem~\ref{thm:p02-w02-p02-002-third-figure}, item~1 with
\(E=\varnothing\)).
\end{proof}
\end{annexitem}
